\section{Literature Search and Survey Methodology}
\label{sec:method}

\subsection{Review design and scope}

This study adopts a structured narrative review to examine the evolution of legal artificial
intelligence from early computational legal reasoning to contemporary retrieval-augmented
generation. The review is organized around the development of methods, the problems they were
designed to address, and the limitations that remain unresolved. This approach is appropriate for
a field in which research spans several related but distinct areas, including legal knowledge
representation, natural language processing, information retrieval, large language models, and
evidence verification.

The purpose is not simply to count publications or rank models by reported performance. Instead,
the review seeks to explain how changes in system design have influenced the way legal information
is represented, retrieved, interpreted, and presented to users. A chronological perspective
establishes the historical progression, while a thematic analysis enables comparison across
studies that address similar technical challenges using different methods.

The review also distinguishes between broad surveys of LegalAI and primary research on particular
architectures or evaluation methods. Broad surveys provide context for understanding the field,
whereas primary studies provide the technical details and empirical evidence needed to compare
specific approaches. For example, \citet{chi2026} organize contemporary LegalAI research around
data, modelling, and evaluation, including legal classification, retrieval, and generation. This
survey uses that broader landscape as context but concentrates on the relationship between
retrieval architecture, preservation of legal structure, and the reliability of generated answers.

The review is therefore guided by four analytical dimensions:
\begin{itemize}[leftmargin=1.4em, itemsep=2pt, topsep=2pt]
\item \textbf{Historical development:} How legal reasoning systems evolved into data-driven
language models and retrieval-augmented systems.
\item \textbf{Technical architecture:} How legal documents are represented, indexed, retrieved,
ranked, and supplied to a generator.
\item \textbf{Reliability and evaluation:} How studies measure retrieval effectiveness, factual
support, hallucination, citation correctness, and overall answer quality.
\item \textbf{Research gaps:} Which challenges remain insufficiently addressed, particularly in
structure-aware retrieval, jurisdiction-specific legal information, and multilingual access.
\end{itemize}

The scope covers foundational research and contemporary developments relevant to these dimensions.
Earlier publications are included where they establish concepts or methods that continue to
influence modern systems. Recent work receives particular attention because the architecture and
evaluation of legal AI are changing rapidly.

\subsection{Information sources and search strategy}

A reproducible review requires more than collecting papers through general web searches. The
sources consulted, search terms used, date of each search, and process for selecting publications
should be documented so that readers can understand how the literature was assembled.

The search strategy should cover both computer science and legal informatics. Relevant studies may
appear in information retrieval, natural language processing, artificial intelligence, law and
technology, and interdisciplinary legal research venues. Accordingly, the recommended sources and
their purposes are summarized in Table~\ref{tab:sources}.

\begin{table*}[t]
\centering
\caption{Recommended information sources for the literature search.}
\label{tab:sources}
\footnotesize
\setlength{\tabcolsep}{6pt}
\renewcommand{\arraystretch}{1.25}
\begin{tabularx}{\textwidth}{>{\raggedright\arraybackslash}p{4.2cm} >{\raggedright\arraybackslash}X}
\toprule
\textbf{Information source} & \textbf{Primary purpose} \\
\midrule
Scopus & Discover peer-reviewed research across computer science, information systems, and legal technology. \\
Web of Science & Identify relevant journal and conference publications and trace citation relationships. \\
IEEE Xplore & Locate research on information retrieval, AI architectures, and computational methods relevant to legal systems. \\
ACM Digital Library & Identify research on NLP, information retrieval, human--computer interaction, and AI system design. \\
ACL Anthology & Locate primary research on legal NLP, language models, retrieval, and language-related evaluation. \\
Google Scholar & Supplement database searches, identify related work, and locate papers not readily indexed elsewhere. \\
Official publisher and proceedings pages & Verify publication metadata, abstracts, publication status, and available full-text records. \\
arXiv & Identify recent preprints and emerging technical approaches, while distinguishing preprints from peer-reviewed publications. \\
\bottomrule
\end{tabularx}
\end{table*}

\textit{Methodological note:} This is the proposed source inventory for the review protocol, not a
claim that every source has already been searched. The final manuscript must identify which
sources were actually consulted and report the dates of those searches.

Searches should combine terms from four groups: the legal domain, the underlying technology,
retrieval and document structure, and reliability. Combining these groups helps identify both
foundational studies and papers that use different terminology for related problems. Example
queries for each theme are listed in Table~\ref{tab:queries}.

\begin{table*}[t]
\centering
\caption{Search concepts and example queries.}
\label{tab:queries}
\footnotesize
\setlength{\tabcolsep}{6pt}
\renewcommand{\arraystretch}{1.25}
\begin{tabularx}{\textwidth}{>{\raggedright\arraybackslash}p{4.0cm} >{\raggedright\arraybackslash}X}
\toprule
\textbf{Search theme} & \textbf{Example search query} \\
\midrule
Legal AI foundations & \texttt{"legal AI" OR "computational legal reasoning" OR "case-based reasoning"} \\
Legal language models & \texttt{("legal NLP" OR "legal language model" OR LEGAL-BERT) AND (classification OR extraction OR reasoning)} \\
Legal retrieval & \texttt{("legal information retrieval" OR "legal case retrieval") AND (ranking OR retrieval OR reranking)} \\
Retrieval-augmented generation & \texttt{("legal RAG" OR "retrieval-augmented generation") AND (law OR legal OR statute OR judgment)} \\
Structure-aware retrieval & \texttt{("structure-aware retrieval" OR "hierarchical indexing" OR "document structure") AND (legal OR statute OR judgment)} \\
Hybrid retrieval and rank fusion & \texttt{("hybrid retrieval" OR "BM25" OR "dense retrieval" OR "reciprocal rank fusion") AND (legal OR jurisprudence)} \\
Reliability and hallucination & \texttt{("legal hallucination" OR "citation verification" OR "factual correctness") AND ("large language model" OR RAG)} \\
Knowledge graphs and multilingual law & \texttt{("legal knowledge graph" OR "knowledge graph RAG" OR "multilingual legal retrieval") AND (AI OR NLP OR retrieval)} \\
Legal benchmarks and evaluation & \texttt{("legal benchmark" OR LegalBench OR LexGLUE OR "legal RAG benchmark") AND (evaluation OR reliability)} \\
\bottomrule
\end{tabularx}
\end{table*}

The queries above are starting points rather than a substitute for database-specific search
syntax. Boolean operators, quotation marks, field restrictions, and supported wildcard characters
differ between platforms. Each query should therefore be adapted to the source being searched,
while preserving its intended concepts.

The search should proceed in two stages. First, broad queries identify the principal research
areas and their foundational publications. Second, focused queries investigate specific
architectural and evaluation issues, including hierarchical document representation, hybrid
retrieval, rank fusion, citation verification, and multilingual legal access. Backward citation
searching, which examines the references of relevant papers, and forward citation searching, which
identifies later work that cites them, can then reveal important studies missed by keyword
searches.

The search period should be recorded explicitly. For a survey covering the development of legal AI
through 2026, foundational work published before the modern LLM era should be included where
relevant, while the contemporary search should extend to the final search date in 2026.
Publications appearing after that date should not be represented as having been covered.

\subsection{Inclusion and exclusion criteria}

Search results can contain many superficially relevant publications that do not provide evidence
useful for the survey. A consistent selection policy is therefore necessary to prevent the final
discussion from becoming an unstructured collection of loosely related papers.

Studies were considered relevant to the survey's scope when they contributed to at least one of
the following: computational legal reasoning, legal NLP, legal information retrieval, legal
language models, RAG architectures, structure-aware document processing, evidence verification, or
legal AI evaluation. Particular importance was given to publications that described their methods
clearly enough to support technical analysis. The full selection criteria are listed in
Table~\ref{tab:criteria}.

\begin{table*}[t]
\centering
\caption{Study selection criteria.}
\label{tab:criteria}
\footnotesize
\setlength{\tabcolsep}{5pt}
\renewcommand{\arraystretch}{1.3}
\begin{tabularx}{\textwidth}{>{\raggedright\arraybackslash}p{2.7cm} >{\raggedright\arraybackslash}X >{\raggedright\arraybackslash}X}
\toprule
\textbf{Criterion} & \textbf{Include} & \textbf{Exclude} \\
\midrule
Topical relevance & Studies addressing legal reasoning, legal NLP, retrieval, generation, verification, or evaluation. & Studies using AI in an unrelated domain without a meaningful connection to legal information processing. \\
Technical contribution & Research presenting a method, architecture, dataset, benchmark, evaluation, or substantive empirical finding. & Publications that mention legal AI only in passing and provide no relevant technical or analytical contribution. \\
Publication type & Peer-reviewed journal articles and conference papers; relevant books and foundational works; clearly identified preprints where they add substantial value. & Duplicate records, unsupported summaries, and non-scholarly material used as if it were primary research. \\
Evidence availability & Publications with accessible bibliographic records and sufficient abstract, full-text, or methodological information for the intended analysis. & Records for which neither the relevant claims nor the publication details can be adequately verified. \\
Evaluation relevance & Studies reporting useful evaluation procedures, datasets, metrics, or clearly described qualitative analyses. & Performance claims whose methods and evidence cannot be established, unless discussed explicitly as an evidence limitation. \\
Jurisdiction and language & Work from any jurisdiction or language when relevant to the review's research questions. & No jurisdiction or language is excluded solely because it is underrepresented in mainstream research. \\
Publication date & Foundational studies from earlier periods and relevant work available by the declared search cutoff. & Work published after the search cutoff, unless clearly identified as a later addition. \\
\bottomrule
\end{tabularx}
\end{table*}

These criteria should be applied with care. For example, a preprint may contain a technically
important method that has not yet undergone peer review. Excluding every preprint could omit
relevant emerging research, while treating every preprint as equivalent to a validated journal
article would overstate the strength of the evidence. The appropriate response is to retain
relevant preprints when justified and label their publication status accurately.

Similarly, a paper should not be included merely because it reports a high score. The evaluation
setting, dataset, baseline, and metric must be sufficiently clear to determine what the reported
result actually demonstrates. A retrieval benchmark, a legal reasoning benchmark, and a human
assessment of citation correctness measure different capabilities and should not be treated as
interchangeable.

\subsection{Screening and source verification}

The selection process should move from broad discovery to detailed examination. The first stage
removes duplicate records and clearly irrelevant results. The second assesses titles and abstracts
against the inclusion criteria. The third examines the full text of potentially eligible studies to
determine whether they provide sufficient evidence for the survey's technical and analytical
objectives.

For each included publication, bibliographic details should be checked against the publisher's
page, the official conference proceedings, or another authoritative publication record. DOI
metadata should be verified where a DOI exists. For preprints, the version consulted and any
subsequently published version should be distinguished to avoid treating multiple versions of the
same work as separate studies.

The full-text assessment should focus on what each study actually establishes. Its proposed
architecture, dataset, retrieval strategy, evaluation design, reported results, and acknowledged
limitations should be recorded separately. When a paper reports an improvement, the review should
identify the baseline and evaluation conditions rather than reproduce the result without context.

This discipline is especially important when comparing legal RAG systems. Two systems may use the
same underlying language model but differ in corpus construction, chunking, retrieval, prompting,
and evaluation. Their reported results cannot automatically be attributed to a single component.
Likewise, a reduction in hallucination on one test set does not establish that a system is
universally more reliable across jurisdictions, legal domains, or languages.

A PRISMA-style flow diagram can make the selection process transparent by reporting records
identified, duplicates removed, records screened, full texts assessed, and studies included,
together with reasons for exclusion. PRISMA 2020 provides established guidance on reporting these
stages and their corresponding counts. However, the use of a flow diagram does not, by itself,
make a review systematic. The counts must come from the actual search and screening process.
Because the current search totals have not been established here, they should be entered only after
the search log has been completed.

\subsection{Data extraction and literature classification}

After screening, the included studies should be analyzed using a consistent extraction framework.
This is necessary because the literature uses different terminology for systems that may share
components, and similar terms may refer to different technical operations. For example, query-time
rank fusion combines the rankings returned by retrieval methods, whereas a model-training procedure
may use preference rankings to select or optimize training examples. Both involve ranking, but they
serve different purposes.

For each study, the extraction process should capture its research objective, legal domain and
jurisdiction, technical approach, evidence sources, evaluation method, main findings, and
limitations. The extracted information should then be classified into thematic categories that
align with the survey's research questions, as shown in Table~\ref{tab:extract}.

\begin{table*}[tp]
\centering
\caption{Literature extraction and classification framework.}
\label{tab:extract}
\footnotesize
\setlength{\tabcolsep}{5pt}
\renewcommand{\arraystretch}{1.3}
\begin{tabularx}{\textwidth}{>{\raggedright\arraybackslash}p{3.0cm} >{\raggedright\arraybackslash}X >{\raggedright\arraybackslash}X}
\toprule
\textbf{Analytical category} & \textbf{Information to extract} & \textbf{Role in the survey} \\
\midrule
Historical foundations & Reasoning paradigm, representation of legal knowledge, and intended task & Explain the origins of computational legal reasoning and the problems inherited by later systems. \\
Legal language models & Model type, training strategy, legal domain, and downstream tasks & Explain how learned representations changed legal text processing. \\
Retrieval architecture & Lexical or dense retrieval, hybrid search, reranking, rank fusion, and filtering & Compare methods used to identify relevant legal evidence. \\
Document representation & Chunking strategy, section hierarchy, metadata, and structural relationships & Assess whether retrieved passages retain the context needed for legal interpretation. \\
Knowledge representation & Ontologies, knowledge graphs, entities, relationships, and provenance & Examine methods for representing explicit connections among legal authorities and concepts. \\
Generation and orchestration & Generator, prompting, evidence use, agents, and workflow design & Explain how retrieved information is converted into legal answers. \\
Reliability and verification & Hallucination, factual support, citation correctness, and verification procedures & Determine how systems assess whether their answers are supported by evidence. \\
Evaluation & Dataset, jurisdiction, benchmark, metrics, baselines, reviewers, and statistical methods & Establish what reported results demonstrate and where comparisons are valid. \\
Practical constraints & Computational requirements, latency, cost, data access, and reproducibility & Identify barriers to deploying and independently evaluating legal AI systems. \\
Research gaps & Limitations acknowledged by authors and unresolved problems identified through comparison & Develop the future research agenda. \\
\bottomrule
\end{tabularx}
\end{table*}

The extracted evidence is synthesized through both chronological and thematic analysis.
Chronological analysis explains why the field moved from hand-crafted legal reasoning toward
learned language representations and retrieval-augmented generation. Thematic analysis then
compares the strengths and limitations of methods addressing similar problems, even when those
methods appeared in different periods.

The synthesis gives particular attention to the distinction between retrieval quality and answer
reliability. Retrieval quality concerns whether relevant and authoritative evidence is found and
ranked appropriately. Answer reliability additionally depends on whether the generator interprets
that evidence correctly, supports its claims with the cited material, and avoids unsupported
conclusions. A system may perform well on one dimension and poorly on another; the survey
therefore treats these as related but distinct evaluation targets.

\subsection{Comparative analysis and methodological limitations}

The final stage compares representative systems using their documented design choices and
evaluation evidence. The comparison considers the retrieval strategy, preservation of legal
document structure, use of external evidence, generation method, verification mechanisms, and
evaluation design. Quantitative results are reported only with their original experimental context.

This restriction is essential because different studies may use different legal jurisdictions,
corpora, query sets, baselines, and definitions of success. For example, a benchmark measuring
legal classification cannot be compared directly with a human evaluation of unsupported legal
claims. Similarly, Recall@$k$, citation precision, and hallucination rate describe different
properties of a system. Combining them into a single ranking without a justified common protocol
would obscure rather than clarify system performance.

The survey also recognizes limitations in the available evidence. First, database coverage and
indexing practices can cause relevant studies to be missed. Second, recent methods may be available
only as preprints or may lack independent evaluation. Third, the terminology used for
structure-aware retrieval, legal grounding, and agentic systems is not always consistent. Fourth,
published results may be difficult to reproduce when datasets, prompts, source documents, or
evaluation procedures are unavailable. Finally, a literature synthesis can identify patterns and
plausible explanations, but it cannot establish that one architecture is universally superior
without controlled, comparable experiments.

For these reasons, the survey distinguishes three kinds of statements throughout the paper:
findings directly reported by individual studies, interpretations drawn by comparing multiple
studies, and open questions for future research. This distinction helps prevent architectural
descriptions from being mistaken for demonstrated empirical advantages.

The resulting methodology supports a critical narrative of the field rather than a simple catalogue
of publications. It establishes a consistent basis for examining how legal AI has evolved, why
modern RAG systems still face reliability challenges, and which design principles may help connect
legal structure, evidence retrieval, and verifiable generation.
